%==============================================================================
\section{Query Decomposition}
\label{sec:decomposition}
%==============================================================================

The \onehop operator of \Cref{sec:onehop} evaluates a single edge pattern
$n_i \rightarrow e \rightarrow n_j$ obliviously, producing a hop relation
$h_e$
of publicly known size $|E_e|$ in a single oblivious sort. A $k$-edge
graph query chains $k$ such patterns together. Rather than treating the
full query as a monolithic join, we observe that it has two properties a
we can exploit.

\emph{Repeated structure.} Every edge in the query is a node-edge-node
traversal --- the exact pattern \onehop is optimized for. Evaluating the full
query as a single join discards this structure entirely, treating edge and
node tables as interchangeable relations.

\emph{Independent structure.} The hop over the first edge does not depend
on the result of the hop over the last edge. Hops at different parts of the
query can be evaluated simultaneously, with results combined only at the
end.

These two observations motivate \emph{query decomposition}: map every edge
in the pattern to a \onehop evaluation, run independent hops in parallel,
and combine the resulting hop-relations into a final answer. The
combination step requires joining $k$ hop-relations on their shared node
ids, where each of $k$ relations have publicly known
sizes. This is precisely what \wei is designed for: it evaluates acyclic
joins obliviously while hiding all intermediate result sizes. The 
query decomposer's goal is to
utilize \onehop wherever possible to efficiently handle the graph traversal
structure and \wei for the residual combination.

We illustrate each step of the decomposition on a running example: a
graph query over a social graph in which person $p_1$ knows
person $p_2$; $p_2$ created comment $c$ that carries tag $t$; and $p_2$
lives in city $\mathit{ci}$, with predicates pinning $p_1$, $t$, and
$\mathit{ci}$:
\begin{center}
\small\ttfamily
p1 -[knows]-> p2 -[created]-> c -[hasTag]-> t \\
\phantom{p1 -[knows]-> }p2 -[livesIn]-> ci \\
\upshape predicates:\ttfamily\ p1.id=933,\ t.id=7,\ ci.id=42
\end{center}
A \emph{query template} defines a node-edge traversal 
pattern and which properties 
are filtered but leaves the predicate values unspecified. A \textit{query},
issued at runtime, instantiates a template with concrete 
predicate values. In the running example, the query template has five 
node labels and four edge labels. We trace the decomposition through 
this example at each step below.

%------------------------------------------------------------------------------
\subsection{Query Templates and Hops}
\label{sec:decomposition:model}
%------------------------------------------------------------------------------

We model a graph query template as a directed \emph{pattern graph} $G_Q = (N,
\mathcal{E})$. Each node $n \in N$ has a label drawn from the schema 
(Section~\ref{sec:background:graph}), and each directed edge 
$e \in \mathcal{E}$ is identified by the triple $(S_e, E_e, D_e)$ --- 
the label of its source node, the edge label, and the label of its 
destination node. A query
also carries a set $P$ of predicates on node or edge properties.

A \emph{hop} $h_e$ is the \onehop evaluation of edge $e$ over its triple
$(S_e, E_e, D_e)$. It has exactly $|E_e|$ rows and carries the edge
attributes of $E_e$ together with the source- and destination-node
attributes, disambiguated by role: in $h_e$, a column is written as
$\mathit{label}\_\texttt{src}\_\mathit{col}$ or
$\mathit{label}\_\texttt{dst}\_\mathit{col}$. This naming lets two hops
that meet at a shared pattern node be joined unambiguously, where the shared node
appears as the destination of one hop and the source of the next, and the
column names make the equality explicit.

\noindent\textbf{Running example.} Each of the four edges becomes one hop:
\begin{center}
\small
\begin{tabular}{@{}ll@{}}
$h_1$: & \texttt{(p1:Person)-[knows]->(p2:Person)} \\
$h_2$: & \texttt{(p2:Person)-[created]->(c:Comment)} \\
$h_3$: & \texttt{(p2:Person)-[livesIn]->(ci:City)} \\
$h_4$: & \texttt{(c:Comment)-[hasTag]->(t:Tag)} \\
\end{tabular}
\end{center}
Note that $h_1$, $h_2$, and $h_3$
are independent of one another --- none consumes another's output --- and
can be evaluated in parallel. Each \onehop reads
only the base node or edge tables and not the output of any other hop.

%------------------------------------------------------------------------------
\subsection{The Decomposition Algorithm}
\label{sec:decomposition:algorithm}
%------------------------------------------------------------------------------

\Cref{alg:decompose} formalizes the procedure in three steps:
(\emph{i})~identify entry nodes and emit one hop per edge via BFS;
(\emph{ii})~stitch hops together on shared nodes; and
(\emph{iii})~place predicates on the appropriate hops. We describe each
step in turn.

\begin{algorithm*}[t]
\caption{Pattern Decomposition}
\label{alg:decompose}
\begin{algorithmic}[1]
\Require query graph $G_Q=(N,\mathcal{E})$ with label-typed nodes and
         directed edges $e=(\mathit{src}(e)\!\rightarrow\!\mathit{dst}(e))$
         typed $(S_e,E_e,D_e)$; node predicates $P$
\Ensure  Hop list $\mathcal{H}$ and reduced query $Q'$
\Statex
\State $\mathit{Roots} \gets \{\, n \in N : n=\mathit{src}(e)\text{ for
       some }e \text{ and } n\neq\mathit{dst}(e)\ \text{for all } e \,\}$
\State $\mathcal{H} \gets [\,]$;\quad $\mathit{seen} \gets \emptyset$
\For{each $r \in \mathit{Roots}$ in fixed order}
  \State $F \gets [\,r\,]$
  \While{$F \neq [\,]$}
    \State $u \gets \textsc{Pop}(F)$
    \For{each out-edge $e=(u\!\rightarrow\!w)$ with $e \notin
         \mathit{seen}$, in fixed order}
      \State $\mathit{seen} \gets \mathit{seen}\cup\{e\}$
      \State append hop $h_e=\langle S_e, E_e, D_e\rangle$ to $\mathcal{H}$
      \State $\textsc{Push}(F, w)$
    \EndFor
  \EndWhile
\EndFor
\Statex
\State $\mathcal{J} \gets \emptyset$
\For{each pattern node $n$ appearing in $\geq 2$ hops}
  \State let $h_{1},\dots,h_{m}$ be those hops in hop order, and
         $\rho_j\in\{\texttt{src},\texttt{dst}\}$ the role of $n$ in $h_j$
  \For{$j=1$ to $m-1$}
    \State $\mathcal{J} \gets \mathcal{J} \cup
           \{\, h_j.(\rho_j\text{-}\mathit{id}) =
           h_{j+1}.(\rho_{j+1}\text{-}\mathit{id}) \,\}$
  \EndFor
\EndFor
\Statex
\For{each predicate $p \in P$ on node $n$}
  \State $h \gets$ first hop in $\mathcal{H}$ with
         $n\in\{\mathit{src}(h),\mathit{dst}(h)\}$
  \State rewrite $p$'s column to $\mathit{label}\_\rho\_\mathit{col}$
         for $n$'s role $\rho$ in $h$
\EndFor
\State $Q' \gets$ join over $\{h\in\mathcal{H}\}$ with conditions
       $\mathcal{J}$ and the placed predicates
\State \Return $\mathcal{H},\ Q'$
\end{algorithmic}
\end{algorithm*}

\emph{Step 1: Entry nodes and hop emission.}
The traversal starts from \emph{entry nodes}: nodes in the query that are 
the source of some edge but the destination of none. A template may
have several entry nodes when it contains branches or disconnected
components; each seeds an independent BFS traversal. From each entry node,
BFS emits one hop per edge in a fixed order, visiting every out-edge of
each dequeued node. The decomposition is exhaustive: every edge becomes a
hop. Emitting hops uniformly keeps the rewrite a function of query
structure alone with no data-dependent choices, and hence no new leakage.

\noindent\textbf{Running example.} Node $p_1$ is the only entry node. BFS
visits $p_2$ next, which has two outgoing edges (\texttt{created} and
\texttt{livesIn}), emitting $h_2$ and $h_3$ in fixed order. From $c$, BFS
emits $h_4$. The four hops are produced in order $h_1, h_2, h_3, h_4$.

\emph{Step 2: Stitching hops.}
A shared node in the pattern graph --- one appearing in more than one hop --- must be
joined on its identity in the residual query. For each such node we collect
the hops in which it appears, note the role it plays in each (source or
destination), and add an id-equality between consecutive occurrences. Along
a chain the shared node is the destination of the earlier hop and the
source of the later, giving a $\mathit{dst}\text{-}\mathit{id} =
\mathit{src}\text{-}\mathit{id}$ condition. At a branch where the shared node is
the source of multiple hops, we get $\mathit{src}\text{-}\mathit{id} =
\mathit{src}\text{-}\mathit{id}$ conditions. Together these reconstruct
exactly the equalities of the original join.

\noindent\textbf{Running example.} Node $p_2$ is a branch: destination of
$h_1$, source of both $h_2$ and $h_3$. Node $c$ is shared between $h_2$
(destination) and $h_4$ (source). Stitching yields three join conditions:
\begin{center}
\small
\begin{tabular}{@{}ll@{}}
\texttt{h1.Person\_dst\_id = h2.Person\_src\_id} & (chain: $p_2$ in
$h_1{\to}h_2$)\\
\texttt{h2.Person\_src\_id = h3.Person\_src\_id} & (branch: $p_2$ in
$h_2, h_3$)\\
\texttt{h2.Comment\_dst\_id = h4.Comment\_src\_id} & (chain: $c$ in
$h_2{\to}h_4$)\\
\end{tabular}
\end{center}

\paragraph{Step 3: Placing predicates.}
Each node predicate is attached to the first hop containing its node and
rewritten against that hop's columns, qualified by the role the node plays
there. We evaluate predicates in \onehop whenever possible, which marks 
each row with a boolean indicating if it satisfied the predicates or not. However, as stated in \cref{sec:onehop}, the size of a \onehop operation
is always bound by its edge list size. 

\noindent\textbf{Running example.} Predicate $p_1.\mathit{id}=933$ lands
on $h_1$ (where $p_1$ is the source); $t.\mathit{id}=7$ lands on $h_4$
(where $t$ is the destination); $\mathit{ci}.\mathit{id}=42$ lands on
$h_3$ (where $\mathit{ci}$ is the destination). The complete residual
query handed to \wei is:
\begin{center}
\small
\begin{tabular}{@{}ll@{}}
$h_1 \bowtie h_2 \bowtie h_3 \bowtie h_4$ & \\[4pt]
\texttt{h1.Person\_dst\_id = h2.Person\_src\_id} & (chain: $p_2$ in $h_1{\to}h_2$)\\
\texttt{h2.Person\_src\_id = h3.Person\_src\_id} & (branch: $p_2$ in $h_2, h_3$)\\
\texttt{h2.Comment\_dst\_id = h4.Comment\_src\_id} & (chain: $c$ in $h_2{\to}h_4$)\\[4pt]
\texttt{h1.Person\_src\_id = 933} & (predicate on $p_1$)\\
\texttt{h4.Tag\_dst\_id = 7} & (predicate on $t$)\\
\texttt{h3.City\_dst\_id = 42} & (predicate on $\mathit{ci}$)\\
\end{tabular}
\end{center}
The residual
join over public-sized hop relations then evaluates them obliviously.
A nine-way join over five node labels and four edge labels becomes a
four-way join over four hop relations, each of publicly known size.

%------------------------------------------------------------------------------
\subsection{Parallelism and Comparison with Naive Execution}
\label{sec:decomposition:parallelism}
%------------------------------------------------------------------------------

To see the concrete benefit of query decomposition, consider a $k$-hop chain query evaluated two
ways. Applying \wei directly to the full $(2k{+}1)$-way join requires six
oblivious sorts per parent-child pair over a join tree of depth $2k$,
totaling $12k$ sorts. Moreover, each 
parent-child pair can only be evaluated in sequence in their scheme.
With decomposition, $k$ \onehop executions run in
parallel (one sort each), followed by \wei over the $k$-way residual join
(six sorts per pair over a tree of depth $k{-}1$, totaling $6(k{-}1)$
sequential sorts). The total is $k + 6(k{-}1) = 7k - 6$ sorts --- a reduction from
$12k$ to $7k{-}6$, with the $k$ \onehop sorts running in parallel rather
than sequentially. For $k=3$ this cuts sequential sort depth from 24 to 12,
and the parallel \onehop sorts contributes little to the critical path. The
gap widens with $k$. This enables our algorithm to handle input tables of XX M, whereas the baseline saturates at size YY.\sujaya{@Bish: please confirm these comparisons}

%------------------------------------------------------------------------------
\subsection{Correctness and Data Independence}
\label{sec:decomposition:correctness}
%------------------------------------------------------------------------------

\paragraph{Correctness.}
Each hop $h_e$ equals the three-way join $S_e \bowtie E_e \bowtie D_e$
(\Cref{sec:onehop}). The stitch conditions reintroduce precisely the
foreign-key equalities of the original pattern, and predicates are
reattached unchanged. The reduced query $Q'$ therefore computes the same
relation as the original $(2k{+}1)$-way join.

\paragraph{Data independence.}
Every choice the procedure makes --- entry nodes, traversal order, hop
emission, stitch conditions, and predicate placement --- is a function of
the query graph and predicate columns alone, both public under our threat
model (\Cref{sec:background:threat}). No data value is read. Each hop has
public size $|E_e|$, and feeding public-sized hops to \wei keeps every
intermediate size public. The composition leaks nothing beyond what the
inputs already reveal.

%------------------------------------------------------------------------------
\subsection{Homogeneous Graphs}
\label{sec:decomposition:homogeneous}
%------------------------------------------------------------------------------

When all edges share one triple --- as in the banking workload of
\Cref{sec:evaluation}, where every edge is typed
$(\texttt{Account})\text{-}[\texttt{Txn}]\text{-}(\texttt{Account})$ ---
all hops are instances of a single hop relation, and the residual query
degenerates to a self-join of that relation under the stitch and predicate
conditions. This homogeneous case requires no algorithmic change and can
benefit from computing the \onehop operator only once.


% \section{Query Decomposition}
% \label{sec:decomposition}
% %==============================================================================

% The \onehop operator of \Cref{sec:onehop} evaluates a single edge pattern
% $n_i \rightarrow e \rightarrow n_j$. A general graph query, however, connects
% $k$ edges into an acyclic join graph and therefore corresponds to a
% $(2k{+}1)$-way join over $k{+}1$ node labels and $k$ edge labels. Handing such a
% query in full to the oblivious multi-way join \wei (\Cref{sec:background:joins})
% discards exactly the structure \onehop was built to exploit: every edge in the
% pattern is itself a one-hop join whose result has the publicly known size $|E|$.

% \system therefore \emph{decomposes} the pattern. Each edge is evaluated by a
% \onehop into a \emph{hop relation} of public size; the residual query then joins
% these hop relations rather than the original node and edge labels. Because each
% hop relation already carries its endpoints' attributes, the residual join is
% narrower: one input per edge instead of $2k{+}1$ base labels. Crucially, the rewrite
% is a \emph{structural} transformation: it reads only the query --- labels, the
% edge connectivity, and the columns named by predicates --- and never inspects a
% single data value, so it introduces no new leakage.

% While \wei evaluates any acyclic join graph, the decomposition is more
% specialized. It assumes the node--edge--node shape of a graph traversal: each hop joins a source
% node, an edge, and a destination node, and the edge relation in the middle
% carries the foreign keys to both endpoints --- exactly the relation \onehop
% enriches from either side. A join with no such middle relation, such as a direct
% node-to-node join, has nothing to fill and falls outside the decomposition. In
% the graph setting this costs nothing: traversals are mediated by edges, so two
% nodes meet only through an edge table, never directly. Every graph pattern
% therefore already has the shape the decomposition assumes, and \wei remains the
% general engine behind it.

% %------------------------------------------------------------------------------
% \subsection{query graphs and Hops}
% \label{sec:decomposition:model}
% %------------------------------------------------------------------------------

% We model a pattern query as a directed \emph{query graph} $G_Q = (N,
% \mathcal{E})$. Each node $n \in N$ carries a label drawn from the schema, and
% each edge $e \in \mathcal{E}$ is directed from a source node $\mathit{src}(e)$ to
% a destination node $\mathit{dst}(e)$. An edge is \emph{typed} by the triple
% $(S_e, E_e, D_e)$ --- its source label, edge label, and destination label. A
% query also carries a set $P$ of predicates on node properties.

% A \emph{hop} $h$ is the \onehop evaluation of one edge $e$ over its triple
% $(S_e, E_e, D_e)$. Following \Cref{sec:onehop}, $h$ has exactly $|E_e|$ rows and
% schema consisting of the edge attributes of $E_e$ together with the source- and
% destination-node attributes, the latter disambiguated by the role in which the
% node appears. We write a hop column as $\mathit{label}\_\texttt{src}\_\mathit{col}$
% or $\mathit{label}\_\texttt{dst}\_\mathit{col}$, recording the node label, the
% role (source or destination), and the original property name. This naming is what
% lets two hops that meet at a shared pattern node be joined unambiguously: the same
% node is reached as the destination of one hop and the source of the next.

% Two facts about a hop are central to the rest of the section. First, a hop's size
% $|E_e|$ is public, so a hop relation can be fed to \wei without revealing any
% selectivity. Second, a hop is type-specific: hops for differently typed edges are
% distinct relations, while hops for identically typed edges share one relation.

% %------------------------------------------------------------------------------
% \subsection{The Decomposition Algorithm}
% \label{sec:decomposition:algorithm}
% %------------------------------------------------------------------------------

% \Cref{alg:decompose} states the procedure. It traverses the query graph,
% emits one hop per edge, reconnects the hops on their shared nodes, attaches the
% predicates, and returns a reduced query $Q'$ for \wei.

% \begin{algorithm*}[t]
% \caption{Pattern Decomposition}
% \label{alg:decompose}
% \begin{algorithmic}[1]
% \Require query graph $G_Q=(N,\mathcal{E})$ with label-typed nodes and directed
%          edges $e=(\mathit{src}(e)\!\rightarrow\!\mathit{dst}(e))$ typed
%          $(S_e,E_e,D_e)$; node predicates $P$
% \Ensure  Hop list $\mathcal{H}$ and reduced query $Q'$
% \Statex
% \State $\mathit{Roots} \gets \{\, n \in N : n=\mathit{src}(e)\text{ for some }e
%        \text{ and } n\neq\mathit{dst}(e)\ \text{for all } e \,\}$
% \State $\mathcal{H} \gets [\,]$;\quad $\mathit{seen} \gets \emptyset$
% \For{each $r \in \mathit{Roots}$ in fixed order} \Comment{BFS: one hop per edge}
%   \State $F \gets [\,r\,]$ \Comment{frontier queue}
%   \While{$F \neq [\,]$}
%     \State $u \gets \textsc{Pop}(F)$
%     \For{each out-edge $e=(u\!\rightarrow\!w)$ with $e \notin \mathit{seen}$, in fixed order}
%       \State $\mathit{seen} \gets \mathit{seen}\cup\{e\}$
%       \State append hop $h_e=\langle S_e, E_e, D_e\rangle$ to $\mathcal{H}$
%       \State $\textsc{Push}(F, w)$
%     \EndFor
%   \EndWhile
% \EndFor
% \Statex
% \State $\mathcal{J} \gets \emptyset$ \Comment{stitch: shared node $\Rightarrow$ id-equality}
% \For{each pattern node $n$ appearing in $\geq 2$ hops}
%   \State let $h_{1},\dots,h_{m}$ be those hops in hop order, and
%          $\rho_j\in\{\texttt{src},\texttt{dst}\}$ the role of $n$ in $h_{j}$
%   \For{$j=1$ to $m-1$}
%     \State $\mathcal{J} \gets \mathcal{J} \cup
%            \{\, h_{j}.(\rho_j\text{-}\mathit{id}) = h_{j+1}.(\rho_{j+1}\text{-}\mathit{id}) \,\}$
%   \EndFor
% \EndFor
% \Statex
% \For{each predicate $p \in P$ on a node $n$} \Comment{place on first hop holding $n$}
%   \State $h \gets$ first hop in $\mathcal{H}$ with $n\in\{\mathit{src},\mathit{dst}\}$
%   \State rewrite $p$'s column to $\mathit{label}\_\rho\_\mathit{col}$ for $n$'s role $\rho$ in $h$
% \EndFor
% \State $Q' \gets$ join over $\{h\in\mathcal{H}\}$ with conditions $\mathcal{J}$ and the placed predicates
% \State \Return $\mathcal{H},\ Q'$
% \end{algorithmic}
% \end{algorithm*}

% \paragraph{Entry nodes.}
% The traversal starts from the pattern's \emph{entry nodes}: those that are the
% source of some edge but the destination of none. These are the heads of the
% pattern, fixed entirely by its connectivity. A pattern may have several entry
% nodes (disconnected components or several in-bound-free starts); each seeds an
% independent traversal.

% \paragraph{One hop per edge.}
% From each entry node we perform a breadth-first traversal that emits one hop for
% every edge, in a fixed order. The decomposition is therefore \emph{exhaustive}:
% every edge of the pattern becomes a hop, with no cost-based selection among
% edges. This is deliberate. Each hop replaces a three-way slice of the join by a
% single $|E|$-sized relation, and emitting them uniformly keeps the rewrite --- and
% hence the access pattern of producing it --- a function of the query alone.

% \paragraph{Branches.}
% A node with more than one outgoing edge yields more than one hop sharing that
% node as their source; the traversal handles such branches with no special case,
% since it simply visits every out-edge of each dequeued node. Chains are the
% special case in which every node has at most one outgoing edge.

% \paragraph{Stitching hops.}
% Wherever two hops meet at the same pattern node, the reduced query must join them
% on that node's identity. We collect, for each node shared by several hops, the
% hops in which it appears together with its role in each, and chain consecutive
% occurrences by an id-equality. Along a path the shared node is the destination of
% the earlier hop and the source of the later one, giving a
% $\mathit{dst}\text{-}\mathit{id} = \mathit{src}\text{-}\mathit{id}$ condition; at a
% branch the node is the source of both hops, giving a
% $\mathit{src}\text{-}\mathit{id} = \mathit{src}\text{-}\mathit{id}$ condition.
% Together these reconstruct exactly the equalities of the original join.

% \paragraph{Placing predicates.}
% Each node predicate is attached to the first hop that contains its node and
% rewritten against that hop's columns, qualified by the role (source or
% destination) the node plays there. Predicates become selection conditions of the
% reduced query, evaluated by \wei rather than inside the operator: as noted in
% \Cref{sec:onehop}, keeping predicates out of the per-hop \onehop preserves its
% $|E|$-sized, selectivity-independent output, so no intermediate size depends on
% how many rows match.

% %------------------------------------------------------------------------------
% \subsection{Example}
% \label{sec:decomposition:example}
% %------------------------------------------------------------------------------

% Consider a heterogeneous pattern over a social graph: a person $p_1$ who knows
% $p_2$; $p_2$ authored a comment $c$ that carries tag $t$; and $p_2$ lives in city
% $\mathit{ci}$, with predicates pinning $p_1$, $t$, and $\mathit{ci}$:
% \begin{center}
% \small\ttfamily
% p1 -[knows]-> p2 -[created]-> c -[hasTag]-> t \\
% \phantom{p1 -[knows]-> }p2 -[livesIn]-> ci \\
% \upshape predicates:\ttfamily\ p1.id=933,\ t.id=7,\ ci.id=42
% \end{center}

% The only entry node is $p_1$ (every other node is some edge's destination). The
% breadth-first traversal emits four hops, each with a \emph{different} triple:
% \begin{center}
% \small
% \begin{tabular}{@{}l@{}}
% $h_1$: \texttt{(Person)-[knows]->(Person)}\\
% $h_2$: \texttt{(Person)-[created]->(Comment)}\\
% $h_3$: \texttt{(Person)-[livesIn]->(City)}\\
% $h_4$: \texttt{(Comment)-[hasTag]->(Tag)}\\
% \end{tabular}
% \end{center}
% Node $p_2$ is a branch: it is the destination of $h_1$ and the source of both
% $h_2$ and $h_3$. Stitching on the shared nodes ($p_2$ and $c$) and placing the
% three predicates on the first hop holding each node yields the reduced query:
% \begin{center}
% \small
% \begin{tabular}{@{}l@{}}
% \texttt{h1.Person\_dst\_id = h2.Person\_src\_id} \quad\textrm{(stitch $p_2$)}\\
% \texttt{h2.Person\_src\_id = h3.Person\_src\_id} \quad\textrm{(branch $p_2$)}\\
% \texttt{h2.Comment\_dst\_id = h4.Comment\_src\_id} \quad\textrm{(stitch $c$)}\\
% \texttt{h1.Person\_src\_id = 933}\\
% \texttt{h3.City\_dst\_id = 42}\\
% \texttt{h4.Tag\_dst\_id = 7}\\
% \end{tabular}
% \end{center}
% A nine-way join over five node labels and four edge labels becomes a four-way
% join over four typed hop relations, each of public size, which \wei evaluates.

% %------------------------------------------------------------------------------
% \subsection{Correctness and Data Independence}
% \label{sec:decomposition:correctness}
% %------------------------------------------------------------------------------

% \paragraph{Correctness.}
% Each hop equals the three-way join $S_e \bowtie E_e \bowtie D_e$ of its edge
% (\Cref{sec:onehop}). The stitch conditions reintroduce precisely the
% foreign-key equalities of the original pattern --- a shared pattern node induces
% the same id-equality whether expressed over base labels or over the hops that
% meet at it --- and the predicates are reattached unchanged. Hence the reduced
% query $Q'$ computes the same relation as the original $(2k{+}1)$-way join.

% \paragraph{Data independence.}
% Every choice the procedure makes --- the entry nodes, the traversal order, which
% edges become hops, the stitch conditions, and where predicates land --- is a
% function of the query graph and the predicate columns, all of which are public
% under our threat model (\Cref{sec:background:threat}). No data value is read.
% The decomposition is thus fixed once the query is fixed, and each hop it produces
% has the public size $|E_e|$; feeding hops of public size to \wei keeps every
% intermediate size public, so the composition leaks nothing beyond the inputs.
% %------------------------------------------------------------------------------
% \subsection{Homogeneous Graphs}
% \label{sec:decomposition:homogeneous}
% %------------------------------------------------------------------------------

% When a query ranges over a single node label and a single edge label --- as in
% the banking workload of \Cref{sec:evaluation}, where every edge is typed
% $(\texttt{Account})\text{-}[\texttt{Txn}]\text{-}(\texttt{Account})$ --- all hops
% share one triple and are therefore instances of a \emph{single} hop relation.
% The reduced query then degenerates to a self-join of that relation under the
% stitch and predicate conditions of \Cref{alg:decompose}. This homogeneous case
% requires no change to the algorithm; it is simply the point at which the typed
% hop relations of the general procedure coincide, and it is the configuration we
% evaluate in \Cref{sec:evaluation}.

% \paragraph{Parallelism.}
% The hops are independent of one another. Each is produced by a separate \onehop
% that reads only its own edge's tables and writes only its own $|E_e|$-sized
% relation; no hop consumes another's result. All hops of a query can therefore be
% computed in parallel, leaving only the residual join over them to run
% sequentially through \wei. Since the number of hops and their sizes are fixed by
% the query, this parallelism is itself data-independent and adds no leakage.
